\chapter{理论基础与关键技术}

\section{大气散射模型}

雾、霾等低能见度现象对成像质量的影响，本质上来源于光在大气中的传播过程受到悬浮粒子的散射与吸收作用。经典的大气散射模型通常表示为：

\[
I(x)=J(x)\cdot t(x)+A\cdot (1-t(x))
\]

其中，\(I(x)\) 表示观测到的退化图像，\(J(x)\) 表示场景真实辐射亮度，\(A\) 表示大气光，\(t(x)\) 为透射率。若进一步考虑场景深度 \(d(x)\) 与散射系数 \(\beta\)，则有：

\[
t(x)=e^{-\beta d(x)}
\]

该模型揭示了一个关键事实：场景越远、雾越浓，透射率越低，图像就越容易“发白”、对比度下降、细节被遮蔽。本文正是利用这一物理启发，将单目深度估计结果与随机采样的 beta 结合，在训练阶段动态生成雾图，并把用于合成的 beta 反向作为监督信号供回归头学习。

需要指出的是，真实高速公路团雾并不一定完全满足理想的均匀散射假设。团雾常常表现为局地浓淡不均、横向分布变化强、边界模糊等特点。因此，若仅按全局均匀透射率进行模拟，合成结果会缺乏足够的空间变化。为解决这一问题，本文在 \texttt{patchy} 模式中进一步引入低频噪声，使有效深度在空间上发生调制，从而得到更符合团雾局地变化特征的合成样本。

\section{单目深度估计}

在没有双目摄像机或激光雷达的条件下，单目深度估计成为恢复场景相对远近信息的重要手段。近年来，MiDaS 系列模型在跨数据集泛化的相对深度估计方面表现较好，能够在不依赖严格度量标定的情况下，为复杂自然图像给出较稳定的深度排序结果。本文采用 MiDaS\_small 作为深度估计器，主要出于两方面考虑：一是其推理开销相对适中，便于对大规模交通视频帧进行离线预计算；二是本文所需并非绝对深度，而是能够支持“远处更易受雾影响”这一基本规律的相对深度先验。

在实现上，本文先对清晰交通图像离线运行 MiDaS\_small，得到归一化深度图，再以 \texttt{序列名\_图像名.npy} 的格式缓存到磁盘。这样在正式训练时，数据集可以直接读取深度缓存，而无需为每个 batch 重复执行深度推理。该设计显著降低了训练期间的算力消耗，并保证了训练、验证阶段对深度信息的稳定访问。

\section{多任务学习}

多任务学习的核心思想在于：若多个任务在输入、场景语义或底层视觉结构上具有相关性，则可以通过共享特征提取层来提升模型效率，并在一定程度上促进不同任务之间的信息互补。对于本文问题而言，车辆检测、雾类型分类和雾浓度回归具有天然联系。雾天气会影响车辆边缘、纹理和对比度表现；车辆目标的可见程度和尺度分布又在一定程度上反映场景能见度水平；高层语义特征中同时包含场景布局、光照退化和目标结构信息，适合承担共享表征角色。

本文采用“共享 YOLO 检测主干 + 分类头 + 回归头”的结构，三项任务联合优化，其总损失可写为：

\[
L = w_{det}L_{det}+w_{cls}L_{fog\_cls}+w_{reg}L_{fog\_reg}
\]

其中，\(L_{det}\) 为目标检测损失，\(L_{fog\_cls}\) 为雾分类交叉熵损失，\(L_{fog\_reg}\) 为 beta 回归均方误差损失。当前实现中三项权重均设置为 1.0，后续也可以根据任务收敛速度和梯度尺度进行进一步平衡。

\section{目标检测与 YOLO 骨干}

YOLO 系列算法由于在检测精度与推理速度之间取得了较好平衡，已广泛用于工业视觉、交通监控和边缘设备部署场景。本文选择 \texttt{yolo11n.pt} 作为基础检测骨干，一方面是因为其参数量和算力开销相对可控，适合毕业设计与原型实现；另一方面，YOLO 主干与颈部已在大规模目标检测任务上具备成熟的多尺度特征提取能力，适合作为共享视觉编码器。

考虑到本文的交通监控场景并不需要完整保留 COCO 的 80 类目标，而主要关注车辆这一粗粒度类别，本文在实现中将所有交通目标统一映射为单类 \texttt{vehicle}，检测头类别数设置为 1。若加载的预训练 YOLO 权重类别数与当前任务不一致，则通过重建单类检测头并保留形状兼容的参数来完成迁移。这种做法兼顾了预训练主干复用与任务定制化需求。

\section{时间平滑与动态阈值}

在视频监控场景中，雾浓度不应当被视为完全独立的逐帧噪声，而更应当具有一定时间连续性。为了避免单帧 beta 回归波动过大影响显示与下游决策，本文在推理阶段采用指数滑动平均（EMA）进行平滑：

\[
\hat{\beta}_{t}=\alpha \beta_t+(1-\alpha)\hat{\beta}_{t-1}
\]

其中 \(\alpha\) 为平滑系数。当前实现中，\texttt{EMA\_ALPHA = 0.1}。随后系统根据平滑后的 beta 值动态调整车辆显示阈值，使系统在低能见度场景下保留更多潜在目标，从而增强监测可见性。这体现了本文方法的一项重要特征：环境感知结果并非孤立输出，而是反过来影响检测展示策略。
